\documentclass[11pt,a4paper]{article}
\usepackage[utf8]{inputenc}
\usepackage[margin=2.2cm]{geometry}
\usepackage{graphicx,booktabs,longtable,array,listings,xcolor}
\usepackage[hidelinks]{hyperref}
\newcommand{\code}[1]{\path{#1}}
\newcommand{\todo}[1]{\textcolor{red}{[#1]}}
% Include an image if the file exists, otherwise print a red note instead of failing.
\newcommand{\maybeimg}[2]{\IfFileExists{#1}{\includegraphics[width=#2]{#1}}{\todo{missing image: \detokenize{#1}}}}
\lstset{basicstyle=\ttfamily\footnotesize,frame=single,breaklines=true,columns=fullflexible}
\setlength{\parindent}{0pt}\setlength{\parskip}{4pt}

\begin{document}

\begin{center}
{\LARGE\bfseries 23CSE455 Design Patterns\\Case Study Report}\\[6pt]
\end{center}

\begin{tabular}{@{}ll@{}}
\textbf{Team:} & B14\\
\textbf{Case study:} & Design Pattern Identification and Evolution in a Next.js-Based\\
 & Digital News and Publishing Application\\
\textbf{Application:} & Next.js Blog CMS (faculty approval: \todo{name, date})\\
\textbf{Date:} & 2026-10-10\\
\textbf{Submission repository:} & 23CSE455\_DP\_CaseStudy
\end{tabular}

\subsection*{Students}
\begin{tabular}{@{}lll@{}}
\toprule
Roll number & Student name & Batch\\
\midrule
AM.SC.U4AIE23121 & Keziah Suzaine George & AIE\\
AM.SC.U4AIE23139 & Nivedita Ramesh & AIE\\
AM.SC.U4CSE23360 & Soorya Krishna & CSE\\
AM.SC.U4CSE23309 & B R Govind & CSE\\
\bottomrule
\end{tabular}

%------------------------------------------------------------
\section{Application and architecture}

\textbf{Original source.}
Repository: \url{https://github.com/thepracticalengineer/nextjs-blog-cms.git}.
Baseline commit: \code{90ec6565877053ee0e93e9f09dcc7c7213aa4832}
(merge of pull request \#102, 8 October 2026). Access date: 2026-10-10. License: MIT.

\textbf{Application.} A full-stack blog CMS with a TipTap editor, Supabase
authentication including MFA, Admin and Author roles, a draft and publication
workflow, comments, tags and categories, public author profiles, newsletter
subscriptions, an AI writing assistant (Claude, Gemini, OpenAI) and a developer REST API.

\textbf{Environment.}
\begin{tabular}{@{}ll@{}}
Node.js & 24.14.0 (repository requests 24.21.0; mismatch noted, nothing failed)\\
Package manager & pnpm 10.29.3\\
Framework & Next.js 16.3.8 (App Router), React 19.3.0, TypeScript 6.0.3\\
Database and auth & Supabase (supabase-js 2.117.2, @supabase/ssr 0.12.7)\\
Other & Zustand 5.0.15, Zod 4.6.5, Vitest 5.0.3, Playwright 1.63.0\\
\end{tabular}

\textbf{Scope.} Inspected: \code{app/}, \code{components/}, \code{features/},
\code{lib/} and \code{proxy.ts}. Excluded: \code{node_modules}, \code{.next},
\code{.git}, \code{.husky}, \code{.github}, \code{.agents}, \code{.claude},
\code{docs/}, \code{public/}, \code{database/}, \code{supabase/migrations/},
\code{e2e/}, \code{__tests__/}, lockfiles and generated files. Reason: not
application logic or not relevant to pattern discovery. Unit tests were used only
as supporting evidence. The source copy in \code{code/} equals the baseline commit
minus \code{.husky}, \code{.github}, \code{tsconfig.tsbuildinfo} and
\code{package-lock.json}.

\textbf{Architecture.}
\begin{figure}[htbp]
\centering
\maybeimg{Part1/architecture.png}{\linewidth}
\caption{Architecture of the application at the baseline commit.}
\end{figure}

The browser reaches public pages, dashboard pages, authentication pages and API
routes. Dashboard and authentication paths pass through \code{proxy.ts} first; its
\code{matcher} covers only \code{/dashboard}, \code{/login}, \code{/register},
\code{/mfa}, \code{/forgot-password} and \code{/reset-password}, so public pages do
not. Pages and routes call domain modules in \code{features/}, which use the
Supabase client creators in \code{lib/supabase}, the permission table in
\code{lib/permissions}, the LLM service and the notification modules.

\textbf{Request flow (admin opens the Users page).} The browser requests
\code{/dashboard/admin/users}. Next.js first runs \code{proxy.ts}, which refreshes
the Supabase session with \code{supabase.auth.getUser()} (line 34), redirects
unauthenticated users to \code{/login}, sends users with an enrolled but unverified
second factor to \code{/mfa}, and, for paths under \code{/dashboard/admin}, reads
the user's role from the \code{profiles} table and redirects non-admins to
\code{/dashboard}. When the request passes, the server component \code{UsersPage}
(\code{app/(dashboard)/dashboard/admin/users/page.tsx}) runs. It calls
\code{requirePermission('users:read')} (line 9), loads data with
\code{getAllUsers()} from \code{features/users/queries} (line 11) and renders
\code{UserTable} (line 21). In the browser, \code{AuthProvider} keeps session state
current by observing \code{supabase.auth.onAuthStateChange}.

\textbf{Run instructions.} See \code{README.md}. Results at the baseline:

\begin{tabular}{@{}p{5.6cm}p{9.6cm}@{}}
\toprule
Command & Result\\
\midrule
\code{pnpm install --frozen-lockfile} & Succeeded, 933 packages\\
\code{pnpm run test:run} & 927 of 927 passed in 83 files on default settings (runs of 71.43 s and 80.10 s). Log: \code{test_result/test_result_vitest.log}\\
\code{pnpm exec vitest run --testTimeout=30000 --maxWorkers=2} & 83 test files passed, 927 of 927 tests passed (222.58 s). Screenshot: \code{test_result/vitest_summary.png}\\
\code{pnpm run typecheck} & Passed. Log: \code{test_result/typecheck.log}; screenshot: \code{test_result/typecheck.png}\\
\code{pnpm run dev}, \code{pnpm run test:e2e} & Not run: Docker is not installed, so local Supabase was unavailable\\
\bottomrule
\end{tabular}

Two further default-setting runs under heavy machine load failed with 6 and 12
failures (130 s and 188 s). Every failure was a 5000 ms timeout in a jsdom UI test,
or a follow-on assertion in the same file after a timeout. The repository sits in a
OneDrive-synced folder. With a 30 s timeout and 2 workers the same code passed
927 of 927, so the failures were timing-related and not application defects.

\begin{figure}[htbp]
\centering
\maybeimg{test_result/vitest_summary.png}{0.9\linewidth}\\[4pt]
\maybeimg{test_result/typecheck.png}{0.9\linewidth}
\caption{Passing unit-test run and typecheck at the baseline commit.}
\end{figure}

%------------------------------------------------------------
\section{Part 1: Pattern discovery}

\textbf{Method.} Candidates were found by reading the file listing and then
reading the likely files in full with line numbers (saved in
\code{Part1/evidence/}). Each candidate was checked against three questions: what
problem it solves here, whether every role maps to a real symbol, and whether the
participants interact as the pattern defines. A class name, a file name or an
LLM statement alone was not accepted as evidence. Combined or absent roles are
stated in the records; no classes were invented.

\subsection{Totals at the baseline}
\begin{tabular}{@{}p{5.5cm}p{9.7cm}@{}}
\toprule
Confirmed pattern types & 3 (Strategy, Observer, Role-Based Access Control)\\
Confirmed pattern instances & 3\\
Concrete participant classes & 7, all functions or components rather than classes: \code{streamClaude}, \code{streamOpenAI}, \code{streamGemini}; \code{AuthProvider}; \code{can}, \code{canAccess}, \code{canAny}\\
Unique participant symbols & 15 application symbols, plus 1 from a dependency (\code{supabase.auth.onAuthStateChange})\\
Rejected candidates & 7\\
Uncertain candidates & 1\\
Observation scope & Files listed in \code{Part1/rejected_uncertain.md}, at the baseline commit\\
\bottomrule
\end{tabular}

\subsection{Pattern inventory}
\begin{tabular}{@{}llcc@{}}
\toprule
Pattern type & Instance name & Instances & Concrete classes\\
\midrule
Strategy & LLM Provider Streaming Strategy & 1 & 3\\
Observer & Auth State Observer & 1 & 1\\
RBAC & RBAC Permission Table & 1 & 3\\
\bottomrule
\end{tabular}

\subsection{Records for each pattern instance}

%---------------- Strategy
\subsubsection*{Instance 1: LLM Provider Streaming Strategy}
\textbf{Pattern type and category:} Strategy, GoF behavioral.
\textbf{Purpose and origin:} streams chat replies from Claude, OpenAI or Gemini
through one call; application code.

\begin{longtable}{@{}p{3.2cm}p{5.8cm}p{6.2cm}@{}}
\toprule
Role & Symbol & File and lines\\
\midrule
\endhead
Strategy & Not separately declared (implicit signature \code{(params: SendMessageParams) => AsyncGenerator<string>}) & \code{features/ai-assistant/llmService.ts}:10--16 (params type)\\
ConcreteStrategy & \code{streamClaude} & same file:69--90\\
ConcreteStrategy & \code{streamOpenAI} & same file:94--118\\
ConcreteStrategy & \code{streamGemini} & same file:129--146\\
Context & \code{sendMessage} & same file:150--155\\
\bottomrule
\end{longtable}

\textbf{Evidence.} The three functions share one parameter type and return type.
The context selects one at run time from \code{params.provider} and delegates:
\begin{lstlisting}
export function sendMessage(params: SendMessageParams): AsyncGenerator<string> {
  if (params.provider === 'claude') return streamClaude(params)
  if (params.provider === 'openai') return streamOpenAI(params)
  if (params.provider === 'gemini') return streamGemini(params)
  throw new Error(`Unsupported provider: ${params.provider}`)
}
\end{lstlisting}
\textbf{Limitation.} There is no declared interface, and selection is a hard-coded
\code{if} chain rather than injection or a registry.

\begin{figure}[htbp]\centering
\maybeimg{Part1/uml_strategy.png}{0.85\linewidth}
\caption{UML for the Strategy instance (source: \texttt{Part1/uml\_strategy.mmd}).}
\end{figure}

\textbf{Verification:} Confirmed. Reviewers: team B14. Method: source reading of
the listed ranges; \code{__tests__/ai-assistant/llmService.test.ts} (20 tests)
passes in the baseline run but was not the sole evidence.
Alternative considered: the repeated provider branches in
\code{generateBlogPost}, \code{generateBlogPostHeadless}, \code{generateChatTitle}
and \code{validateProviderKey} (lines 157--386) do not apply Strategy; see the
rejected list.

%---------------- Observer
\subsubsection*{Instance 2: Auth State Observer}
\textbf{Pattern type and category:} Observer, GoF behavioral.
\textbf{Purpose and origin:} keeps UI session and profile state in sync with login,
logout and token refresh; application code using the Supabase dependency.

\begin{longtable}{@{}p{3.2cm}p{5.8cm}p{6.2cm}@{}}
\toprule
Role & Symbol & File and lines\\
\midrule
\endhead
Subject & \code{supabase.auth} (dependency, not application code) & \code{features/auth/context/AuthProvider.tsx}:37--39\\
ConcreteObserver & \code{AuthProvider}, via the callback given to \code{onAuthStateChange} & same file:9--59 (callback 39--48)\\
Detach & \code{subscription.unsubscribe()} & same file:50\\
State holder & \code{AuthContext}, \code{AuthContext.Provider} & \code{AuthContext.tsx}:14--19; \code{AuthProvider.tsx}:54--58\\
Consumers & \code{useAuth}; \code{useSession} & \code{AuthContext.tsx}:21--23; \code{features/auth/hooks/useSession.ts}:5--8\\
\bottomrule
\end{longtable}

\textbf{Evidence.} The component subscribes, updates state when the callback
fires, and unsubscribes in the effect cleanup:
\begin{lstlisting}
const { data: { subscription } } = supabase.auth.onAuthStateChange(
  async (_event, session) => { setSession(session); setUser(session?.user ?? null); ... })
return () => subscription.unsubscribe()
\end{lstlisting}
The subject belongs to \code{@supabase/supabase-js}, so no subject class is
invented. React Context is the delivery mechanism to consumers and is not counted
as a second instance.

\begin{figure}[htbp]\centering
\maybeimg{Part1/uml_observer.png}{0.85\linewidth}
\caption{UML for the Observer instance (source: \texttt{Part1/uml\_observer.mmd}).}
\end{figure}

\textbf{Verification:} Confirmed. Reviewers: team B14. Method: source reading of
\code{AuthProvider.tsx}, \code{AuthContext.tsx} and \code{useSession.ts}.

%---------------- RBAC
\subsubsection*{Instance 3: RBAC Permission Table}
\textbf{Pattern type and category:} Role-Based Access Control, architecture and
security pattern (not a GoF pattern). \textbf{Purpose and origin:} centralizes
which role may perform which action; application code.

\begin{longtable}{@{}p{3.2cm}p{5.8cm}p{6.2cm}@{}}
\toprule
Role & Symbol & File and lines\\
\midrule
\endhead
Permission table & \code{rolePermissions} & \code{lib/permissions/index.ts}:3--30\\
Check & \code{can} & same file:32--35\\
All-of check & \code{canAccess} & same file:37--40\\
Any-of check & \code{canAny} & same file:42--45\\
Types & \code{Role}; \code{Permission} & \code{lib/permissions/types.ts}:1; 3--18\\
\bottomrule
\end{longtable}

\textbf{Evidence.} \code{rolePermissions} maps each of the roles \code{admin} and
\code{author} to a list of permissions. \code{can} looks up the table, and
\code{canAccess} and \code{canAny} apply \code{every} and \code{some} over
\code{can}:
\begin{lstlisting}
export function can(role: Role | null | undefined, permission: Permission): boolean {
  if (!role) return false
  return rolePermissions[role]?.includes(permission) ?? false
}
\end{lstlisting}

\begin{figure}[htbp]\centering
\maybeimg{Part1/uml_rbac.png}{0.8\linewidth}
\caption{UML for the RBAC instance (source: \texttt{Part1/uml\_rbac.mmd}).}
\end{figure}

\textbf{Verification:} Confirmed as an architecture pattern. Reviewers: team B14.
Method: source reading; \code{__tests__/lib/permissions.test.ts} (29 tests) passes
in the baseline run.

\subsection{Rejected and uncertain candidates}
\textbf{Uncertain (1).} \code{proxy.ts} guard pipeline (lines 4--117): sequential
guard checks in one function. It resembles Chain of Responsibility, but there is no
handler interface, no successor links and no separate handlers, so roles cannot be
mapped without inventing classes. Classified as a Next.js framework idiom.

\textbf{Rejected (7).}
\begin{longtable}{@{}cp{5cm}p{9.5cm}@{}}
\toprule
\# & Candidate & Reason\\
\midrule
\endhead
R1 & Supabase client creators as Factory Method & Four standalone functions chosen by import path; no creator hierarchy or polymorphic selection; \code{createStaticClient} can return null, so they share no contract\\
R2 & User-confirmed notifications as Observer & \code{route.ts}:67--74 calls two functions in \code{Promise.all}, each with its own catch; no subscriber list or notifier interface\\
R3 & \code{sendNewsletterEmail} & Single function wrapping Resend; one participant\\
R4 & Publication validation as State or Specification & \code{status} is a plain string; one function with sequential checks, no rule objects\\
R5 & Repeated provider \code{if} chains & Strategy is not applied; branching duplicated (Part 2 refactoring target)\\
R6 & Helper functions & \code{buildSystemPrompt}, \code{resolveGeminiModel}, \code{isRateLimitError}, \code{escapeHtml}, \code{sanitizeSubject}: plain utilities\\
R7 & React Context as a separate Observer & Folded into the Auth State Observer to avoid double counting\\
\bottomrule
\end{longtable}

%------------------------------------------------------------
\section{Part 2: Evolution and refactoring}

\textbf{Scope and traceability.} The refactoring commits currently documented in Part 2 are:
\begin{itemize}
\item CH01 provider inference: branch \code{part2/ch01-provider-refactor}, commit \code{e7dfd1c}; changes \code{features/ai-assistant/llmService.ts} and \code{__tests__/ai-assistant/llmService.test.ts}.
\item CH02 publication validation: branch \code{part2/ch02-publishing-workflow}, commit \code{7be15fb}; changes \code{features/posts/publication.ts}.
\item CH03 shared generated-post JSON parsing: branch \code{part2/ch03-json-parsing-refactor}, commit \code{6ce7157}; changes \code{features/ai-assistant/llmService.ts}.
\end{itemize}
The reproducible commit map is in \code{Part2/code/README.md}. The source snapshot stored in \code{B14/code/} is the Part 1 baseline copy; it does not automatically include the commits on the separate application branches.

\textbf{Scope mismatch requiring team/instructor confirmation.} The initial Part 2 prompt specifies CH01 article recommendations, CH02 content moderation, and CH03 subscription/paywall access control. The three refactoring commits listed above instead centralize provider inference, extract publication validation helpers, and share generated-post JSON parsing. A separate application branch \code{part2/ch01-article-recommendations}, commit \code{68cd127}, implements related-article recommendations. The evidence examined for this report does not establish that content moderation or server-side premium-content access control was implemented. This report therefore describes the confirmed code changes and does not claim that moderation or paywall requirements are complete.

\subsection*{CH01: Centralize AI model-provider inference}
\textbf{Smell and technique.} Repeated model-name inference logic; extract \code{inferProviderFromModel()} and use it in \code{generateBlogPost()}.
\textbf{Evidence.} Commit \code{e7dfd1c} changes the LLM service and its unit tests. See \code{changes/Change1/implementation.md} and focused test logs. The latest focused log committed to this repository shows 20 tests passed; the implementation note reports 23 focused tests and a full-suite total of 1,025, which have not been reconciled with that latest log. These historical figures should not be presented as the latest result without their corresponding original logs.
\textbf{Pattern effect.} The confirmed \emph{LLM Provider Streaming Strategy} is preserved: \code{sendMessage()} and its three streaming functions are unchanged by this commit. The helper centralizes a separate blog-generation inference responsibility; it does not, by itself, establish a new Strategy instance.

\subsection*{CH02: Extract publication validation helpers}
\textbf{Smell and technique.} The publication validation function contained inline title and readable-body checks; extract \code{validatePublicationTitle()} and \code{validatePublicationContent()} in \code{features/posts/publication.ts}.
\textbf{Evidence.} Commit \code{7be15fb}; focused publication test evidence is in \code{changes/Change2/evidence/targeted-tests.log}. The implementation record reports 30 baseline and 30 after-refactoring publication tests, a passing typecheck, and 1,022 tests passing across 88 files. Read the test result alongside the stored raw logs.
\textbf{Pattern effect.} This improves separation of concerns but introduces no claimed GoF pattern. The confirmed Strategy, Observer and RBAC instances are preserved.

\subsection*{CH03: Extract shared generated-post JSON parsing}
\textbf{Smell and technique.} Duplicate JSON cleanup, parsing and invalid-JSON error handling in \code{generateBlogPost()} and \code{generateBlogPostHeadless()}; extract \code{parseGeneratedPostJson()}.
\textbf{Evidence.} Commit \code{6ce7157}; its diff shows the helper and both call sites. The latest focused log shows 20 tests passed. The full-suite total reported in \code{changes/Change3/implementation.md} is a separate recorded run and must not be inferred from the focused log.
\textbf{Pattern effect.} The helper extraction does not introduce a GoF pattern. Strategy, Observer and RBAC remain unchanged.

\subsection*{Pattern evolution and counts}
The source-based rows and status justifications are in \code{Part2/pattern_evolution.csv}; counts use the Part 1 count schema in \code{Part2/Pattern_count.csv}. The three confirmed instances (Strategy, Observer and RBAC) remain one instance each across the three refactoring snapshots. The extracted helpers are not counted as pattern instances: names alone do not establish a pattern.

\subsection*{LOC and measurement limitation}
Commit diffs supply changed-line statistics, but a consistent before/after source LOC measurement has not yet been captured across identical paths and a declared counting rule. This report does not invent an LOC delta. To complete that metric, run the same LOC tool and counting rule against the relevant before and after commits and record the command and output. Full application runtime/end-to-end testing was not run in the baseline environment because local Supabase/Docker was unavailable, as noted in Part 1.

\section{Part 3: Cross-language comparison}
\todo{Not started.}

%------------------------------------------------------------
\section{LLM records and human verification}

\textbf{Run and model.} Part 1 pattern analysis used an LLM chat assistant on
2026-10-10. The displayed model name, exact prompts, raw outputs and the
evaluation of each analysis are recorded in \code{Part1/llm.csv},
\code{Part1/Part1_prompt.txt} and \code{Part1/llm/}.

\textbf{Claims and evaluation.} Every pattern claim was checked against the
numbered source listings in \code{Part1/evidence/}. Corrections made after
verification: the line number given for \code{getUser()} in \code{proxy.ts} was
wrong and the true line is 34; a claim that the Users page uses \code{can()} and
the server Supabase client was unsupported, since the page calls
\code{requirePermission} and \code{getAllUsers}, so it was removed; and an initial
Observer guess for the user-confirmed notification fan-out was rejected after
reading \code{route.ts}. No generated code was used in Part 1.

%------------------------------------------------------------
\section{Findings, contributions and references}

\textbf{Findings.} At the baseline commit, three pattern instances were confirmed
(Strategy, Observer, RBAC), one candidate is uncertain and seven were rejected.
The application mostly relies on framework idioms and plain functions rather than
classical GoF structures. The Strategy is function-based with an \code{if}-chain
selector and no declared interface. The Observer subject comes from a dependency.
RBAC is an architectural pattern. Of the eight areas in the brief, articles,
authors, categories, comments, authentication and newsletter subscription exist;
moderation exists only as comment deletion, and recommendations and a paywall were
not found. Limitations: the app was not run (no Docker), and behavior was verified
from source and the unit-test suite only.

\textbf{Contributions.}
\begin{tabular}{@{}p{3.3cm}p{3.6cm}p{1.2cm}p{6.6cm}@{}}
\toprule
Roll number & Name & Batch & Contribution\\
\midrule
AM.SC.U4AIE23121 & Keziah Suzaine George & AIE & \todo{actual contribution}\\
AM.SC.U4AIE23139 & Nivedita Ramesh & AIE & \todo{actual contribution}\\
AM.SC.U4CSE23360 & Soorya Krishna & CSE & \todo{actual contribution}\\
AM.SC.U4CSE23309 & B R Govind & CSE & \todo{actual contribution}\\
\bottomrule
\end{tabular}

\textbf{References.}
\begin{enumerate}
\item E.~Gamma, R.~Helm, R.~Johnson, J.~Vlissides, \emph{Design Patterns: Elements of Reusable Object-Oriented Software}, Addison-Wesley, 1994.
\item R.~Sandhu, E.~Coyne, H.~Feinstein, C.~Youman, ``Role-Based Access Control Models,'' \emph{IEEE Computer} 29(2), 1996.
\item Next.js Blog CMS, \url{https://github.com/thepracticalengineer/nextjs-blog-cms}, commit 90ec656, MIT license.
\item Next.js documentation, \url{https://nextjs.org/docs}.
\item Supabase documentation, \url{https://supabase.com/docs}.
\item Mermaid, \url{https://mermaid.js.org}.
\end{enumerate}

\section{Team folder and datasets}
\begin{figure}[htbp]\centering
\maybeimg{test_result/github_team_folder.png}{\linewidth}
\caption{Team folder B14 in the course repository.}
\end{figure}
\begin{figure}[htbp]\centering
\maybeimg{test_result/github_part1.png}{\linewidth}
\caption{Part1 folder in the course repository.}
\end{figure}

\end{document}